\section{Methodology}
\label{sec:method}

Our design isolates one variable---post-training intensity---and measures its effect on
internal answer commitment, behavioral reactivity, and uncertainty redistribution on
\emph{identical} items.

\subsection{Models: a matched post-training gradient}
\label{sec:models}

We use three checkpoints from one model lineage, sharing architecture and tokenizer so that
the systematic difference between them is the amount and type of post-training
(\tabref{tab:models}). The \base model is pretrained only; \instruct adds supervised
fine-tuning plus RLHF/DPO; and \rlm is distilled from DeepSeek-R1's RL reasoning traces.
Because true base models are not served by any chat API, all weights are run locally. Each
model is used in its \emph{native} interface: \base receives a plain transcript prompt (it
has no chat template), while \instruct and \rlm use their own chat templates. We run bf16
weights on an NVIDIA RTX A6000, with greedy decoding for commitment measures and temperature
$0.7$ sampling for entropy. Seeds are fixed (42; 123 for Experiment~3 sampling).

\begin{table}[t]
    \centering
    \caption{The matched post-training gradient. All three models share the \qwenbase
    architecture and tokenizer lineage, so post-training stage is the only systematic
    difference.}
    \label{tab:models}
    \resizebox{\textwidth}{!}{%
    \begin{tabular}{@{}llp{6.0cm}@{}}
        \toprule
        \textbf{Key} & \textbf{HuggingFace id} & \textbf{Post-training stage} \\
        \midrule
        \base     & \texttt{Qwen/Qwen2.5-1.5B}                  & none (pretrained only) \\
        \instruct & \texttt{Qwen/Qwen2.5-1.5B-Instruct}         & SFT + RLHF/DPO \\
        \rlm      & \texttt{deepseek-ai/DeepSeek-R1-Distill-Qwen-1.5B} & distilled from R1's RL reasoning traces \\
        \bottomrule
    \end{tabular}%
    }
\end{table}

\subsection{Datasets}
\label{sec:data}

We use three pre-gathered local datasets chosen to fit each experiment. \arc (science
multiple-choice) gives clean single-letter answers and is used for the reactivity and probe
experiments. \gsm (grade-school math) gives numeric answers and serves as a robustness check
for reactivity. \strat (decomposable yes/no questions with an accompanying \texttt{facts}
chain) supports the rotating-certainty experiment, because each item can be split into two
hops.

\subsection{Experiment 1: behavioral reactivity (C2b)}
\label{sec:exp1}

We convert each single-turn item into a two-turn episode: the model answers, feedback is
injected, and the model answers again. Feedback is one of three types. \generic is
content-free pushback (``that's incorrect, try again''), following the ``Try
Again''/UFO protocol; \challenge is the FlipFlop-style ``are you sure?''; and \hint is genuine
information pointing toward the correct answer. We report the \emph{answer change rate} from
turn~1 to turn~2 and, to reveal \emph{how} a model reacts, the directional flip rates
$\Delta^{i\rightarrow c}$ (a wrong answer corrected to right---a beneficial fix) and
$\Delta^{c\rightarrow i}$ (a right answer broken into wrong---harmful over-reaction). We run
$N=200$ \arc items per feedback type and $N=150$ \gsm items as a robustness check. Note that
the \gsm hint is non-informative (``re-check your arithmetic''), so on \gsm all three feedback
types probe reaction to \emph{content-free} nudges; we flag this when interpreting.

\subsection{Experiment 2: commitment-earliness probe (C1, C2a)}
\label{sec:exp2}

We greedily generate each model's answer, parse the model's \emph{own} final letter, and
train a balanced logistic-regression probe (80/20 split) to predict that letter from hidden
states. We probe two locations: the \textbf{prompt-end} hidden state (across all layers),
which measures how much of the answer is pre-encoded \emph{before} generation, and a
mid-layer trajectory \emph{along} the generation timeline, which measures how early the answer
becomes decodable. Probe accuracy at the prompt position far above the four-way chance level
($0.25$) indicates the answer is pre-encoded. Because the probe targets the model's own
answer rather than the gold answer, a model's low task accuracy does not bias the measurement.
We use $N=440$ four-way (A--D) \arc items.

\subsection{Experiment 3: rotating certainty (C3)}
\label{sec:exp3}

For each \strat item we split the \texttt{facts} chain into two hops, $A$ and $B$. We then
compare three conditions: \texttt{baseline} (neither hop given), \texttt{+A} (hop $A$ held
certain in the prompt), and \texttt{+B} (hop $B$ held certain). Our uncertainty measure is the
predictive binary entropy of the yes/no answer, estimated from $K=6$ samples per item, which
is model-agnostic and does not depend on logit access. We report $H$(baseline), $H$(+A),
$H$(+B), the signed reductions $\Delta H$, and paired significance. A reduction in entropy
when a hop is fixed indicates uncertainty was \emph{redistributed} from that hop onto the rest
of the question. We run $N=120$ items.

\subsection{Baselines and statistics}
\label{sec:stats}

The within-study baseline for every comparison is the \base model. Probe accuracy is compared
against the random/majority chance level ($0.25$ for four-way \arc). We report bootstrap 95\%
confidence intervals (10k resamples), Fisher exact tests for the categorical change rates in
Experiment~1, Wilcoxon signed-rank tests for the paired entropy reductions in Experiment~3,
and Mann--Whitney $U$ for cross-model comparisons. Parse-failure rates are logged ($\leq6\%$,
concentrated in the reasoning model) and reported rather than silently dropped.
